\chapter{Pipeline Orchestration with Amazon MWAA and Step Functions}
\index{Apache Airflow}\index{Amazon MWAA}\index{AWS Step Functions}\index{DAG}\index{orchestration}\index{XCom}\index{sensors}%
\begin{objectives}
\begin{itemize}
    \item Apply core Apache Airflow concepts: DAGs, operators, tasks, scheduling, retries, and XComs.
    \item Explain Amazon MWAA architecture and its operational model versus self-managed Airflow.
    \item Use AWS Step Functions for native serverless orchestration and know when it beats Airflow.
    \item Author a production-grade DAG that waits for S3 arrival, triggers a Glue job, runs an Athena validation, and alerts on failure.
    \item Design a resilient, self-healing orchestration layer for a 15-system banking reconciliation.
    \item Assemble the Milestone~3 automated batch ELT pipeline with a data-quality gate.
\end{itemize}
\end{objectives}

\begin{crosscloud}
Orchestration splits into two families everywhere: general-purpose DAG schedulers and native state machines. Both exist on every cloud.

\vspace{2mm}
{\footnotesize
\begin{tabular}{@{}p{2.8cm}p{3.0cm}p{3.0cm}p{3.0cm}@{}}
\toprule
\textbf{Capability} & \textbf{AWS} & \textbf{Azure} & \textbf{GCP} \\
\midrule
Managed Airflow & Amazon MWAA & Data Factory managed Airflow & Cloud Composer \\
Native state machine & Step Functions & Azure Logic Apps / Durable Functions & Workflows \\
File-arrival trigger & EventBridge + sensor & Event Grid trigger & Eventarc trigger \\
Failure alerting & SNS via operator/rule & Monitor alerts & Alerting to Pub/Sub \\
\addlinespace
Inter-task data & XComs (small) / S3 (large) & Pipeline variables / ADLS & XComs / GCS \\
Retry semantics & Task retries + backoff & Activity retry policy & Step retry blocks \\
\bottomrule
\end{tabular}}

\vspace{1mm}
\textbf{Fabric} pipelines play the orchestrator role inside the SaaS boundary. \textbf{Snowflake} tasks and DAGs orchestrate in-warehouse work, and dbt (Chapter~24) orchestrates transformation graphs---the same dependency idea at the SQL layer. \textbf{MongoDB} triggers and Atlas schedulers cover the document-store corner.
\end{crosscloud}

\section{Week 12 Roadmap and Mental Model}
Weeks 9--11 built the engines; Week~12 conducts them. An orchestrator answers four questions a scheduler cron entry cannot: What depends on what? What happens when a step fails? What data does a run produce, and which run produced it? And who gets paged? This week's two instruments approach those questions from different histories: \textbf{Airflow} (via Amazon MWAA) is the general-purpose, code-first DAG platform the industry standardized on \cite{apacheAirflow}; \textbf{Step Functions} is AWS's native, serverless state machine with deep service integrations and per-transition billing \cite{awsStepFunctionsDocs}.

The professional skill is not picking one---it is knowing the boundary. Rule of thumb this chapter defends: Step Functions for event-driven, high-rate, AWS-native task graphs; Airflow for schedule-driven, dependency-heavy data estates with cross-day logic and a Python-native team.

\begin{definitionbox}
\textbf{Orchestration} is the managed execution of a directed acyclic graph (DAG) of data tasks with dependency resolution, retry and alerting policy, run history, and parameterized scheduling---distinct from the \emph{computation} the tasks themselves perform.
\end{definitionbox}

\begin{keyterms}
\textbf{DAG}: the code-defined graph of tasks and dependencies. \textbf{Operator}: a task template (run Glue job, run Athena query). \textbf{Sensor}: a polling operator that waits for a condition (file exists). \textbf{XCom}: Airflow's small inter-task message channel. \textbf{catchup}: whether Airflow backfills missed schedules on deploy. \textbf{State machine}: Step Functions' Amazon States Language definition of task flow. \textbf{Execution role / task role}: MWAA's separation between the environment's identity and per-task AWS permissions.
\end{keyterms}

\section{Monday: Airflow Concepts}
\subsection{DAGs, Operators, Tasks, and Scheduling}
A DAG is Python code declaring tasks and their dependency edges; a \emph{DAG run} is one scheduled or triggered execution; a \emph{task instance} is one task within one run. Operators wrap actions---\texttt{GlueJobOperator}, \texttt{AthenaOperator}, \texttt{SnsPublishOperator}---and sensors poll for external conditions like file arrival. Scheduling is cron or timetable-based, with two settings that bite newcomers: \texttt{catchup} (if true, deploying a DAG with a past \texttt{start\_date} instantly schedules every missed run) and retries with backoff, which convert transient failures---throttling, a brief S3 inconsistency---into eventual success instead of pages \cite{apacheAirflow}.
\index{catchup}\index{retry}\index{sensor}

\begin{pitfall}
\textbf{Deploying with \texttt{catchup=True} and a past \texttt{start\_date}.} The scheduler will cheerfully launch months of backfill runs against production systems the moment the DAG file lands. Set \texttt{catchup=False} unless you have designed the backfill, and design the backfill deliberately when you need one.
\end{pitfall}

\subsection{XComs and the Data Rule}
XComs pass small values between tasks---a row count, an S3 key, a status. They are stored in the metadata database, so they are \emph{not} a data channel: the rule is \textbf{data flows through S3, signals flow through XComs}. A task that XComs a dataframe will bloat the metadata database and fail mysteriously at scale; a task that writes S3 and XComs the key scales indefinitely.

\section{Tuesday: Amazon MWAA Architecture}
\subsection{Managed Airflow, Honest Boundaries}
Amazon MWAA runs Airflow for you: a managed scheduler and web server, with workers autoscaling on ECS Fargate, all inside your VPC \cite{awsMWAADocs}. You choose an environment class, a worker range, and a requirements file. DAGs deploy by syncing an S3 bucket---the \emph{dags} prefix is the deployment surface, so CI/CD for DAGs is ``merge to main, sync to S3'' (Chapter~24 formalizes this).

The operational realities: environment updates cause brief scheduler restarts; worker scaling lags sudden parallelism spikes; and every AWS call from a DAG uses the execution role, so least-privilege means that role can call exactly the Glue jobs, Athena workgroups, and SNS topics the DAGs need---no more.
\index{MWAA execution role}

\begin{tipbox}
Wire MWAA and Glue failures to CloudWatch alarms that publish to an SNS topic subscribed by the on-call runbook. Alarm on \texttt{FailedTasks} and DAG-duration anomalies---never rely on someone noticing a red box in the Airflow UI. Every alert must link to a runbook entry with triage steps.
\end{tipbox}

\section{Wednesday: AWS Step Functions}
\subsection{Native Serverless Orchestration}
Step Functions defines workflows in Amazon States Language (JSON/YAML): Task states call AWS services (with optimized integrations that handle polling and callbacks), Choice states branch on payload fields, Map states fan out, and built-in retry/catch blocks implement per-state failure policy \cite{awsStepFunctionsDocs}. Standard workflows run up to a year with full execution history; Express workflows handle high-rate, sub-five-minute event processing.

Step Functions' signature advantage for data pipelines is the \emph{service integrations}: triggering a Glue job, running an EMR Serverless job, or executing an Athena query is a one-line Task state with IAM scoped per state. Its signature constraint is the 256~KB payload limit between states---the same ``S3 for data, signals for state'' rule as XComs, enforced by the platform.
\index{Step Functions}\index{Amazon States Language}

\subsection{Choosing Between MWAA and Step Functions}
\begin{table}[H]
\centering
\small
\begin{tabular}{@{}p{3.4cm}p{4.4cm}p{4.4cm}@{}}
\toprule
\textbf{Criterion} & \textbf{MWAA (Airflow)} & \textbf{Step Functions} \\
\midrule
Sweet spot & Schedule-driven estates, cross-day dependencies, rich backfill & Event-driven, high-rate, AWS-service-centric flows \\
Programming model & Python DAGs, huge operator ecosystem & ASL JSON/YAML, visual designer \\
Billing & Environment-hours (always on) & Per state transition (per request for Express) \\
Data passing & XComs (discipline required) & Payloads capped at 256 KB (enforced) \\
Team fit & Python-native data engineers & Platform/app teams, event architectures \\
\bottomrule
\end{tabular}
\caption{Choosing the orchestrator: MWAA versus Step Functions.}
\label{tab:ch12-orchestrators}
\end{table}

\section{Thursday Hands-on Project: A Production-Shaped MWAA DAG}
\index{hands-on project!MWAA DAG}
Author and deploy a DAG that waits for a file in S3, triggers a Glue job, validates the load with Athena, and publishes an SNS alert on any failure.

\subsection*{Lab Objectives}
\begin{itemize}
    \item Deploy the DAG to an MWAA environment's S3 DAG bucket.
    \item Configure the execution role with least-privilege access to exactly the four services used.
    \item Prove the failure path: force the Glue job to fail and confirm the SNS alert fires.
    \item Explain every scheduling-related setting in the DAG header.
\end{itemize}

\subsection*{Procedure}
The DAG below is the lab's complete artifact. Read it once for the task graph and once for the operational settings: \texttt{catchup=False}, retries with delay, a tuned sensor, and a failure-triggered alert.

\begin{lstlisting}[language=Python,caption={The daily reconciliation DAG for MWAA.},label={lst:ch12-dag}]
from datetime import datetime
from airflow import DAG
from airflow.providers.amazon.aws.sensors.s3 import S3KeySensor
from airflow.providers.amazon.aws.operators.glue import GlueJobOperator
from airflow.providers.amazon.aws.operators.athena import AthenaOperator
from airflow.providers.amazon.aws.operators.sns import SnsPublishOperator

with DAG(
    dag_id="daily_recon_pipeline",
    schedule="@daily",
    start_date=datetime(2026, 1, 1),
    catchup=False,
    default_args={"retries": 2, "retry_delay": 300},
) as dag:
    wait_for_file = S3KeySensor(
        task_id="wait_for_file",
        bucket_name="bank-raw-ingest",
        bucket_key="recon/{{ ds }}/transactions.csv",
        poke_interval=300, timeout=3600,
    )
    run_glue = GlueJobOperator(
        task_id="run_glue", job_name="recon-etl", iam_role_name="GlueExecRole",
    )
    validate = AthenaOperator(
        task_id="validate", query="SELECT COUNT(*) FROM recon.daily_load",
        database="recon", output_location="s3://bank-athena-results/",
    )
    alert = SnsPublishOperator(
        task_id="alert", target_arn="arn:aws:sns:us-east-1:123456789012:data-alerts",
        message="Recon pipeline failed for {{ ds }}", trigger_rule="one_failed",
    )
    wait_for_file >> run_glue >> validate >> alert
\end{lstlisting}

Deploy by uploading to the environment's DAG prefix:

\begin{lstlisting}[language=bash,caption={Deploying the DAG and watching it appear.},label={lst:ch12-deploy}]
$MwaaEnv = "week12-mwaa"
$DagBucket = (aws mwaa get-environment --name $MwaaEnv `
  --query 'Environment.DagS3Path' --output text) -replace '/dags$',''

aws s3 cp daily_recon_pipeline.py "s3://$DagBucket/dags/"

# The scheduler picks it up within a minute or two:
aws mwaa get-environment --name $MwaaEnv --query 'Environment.Status'
\end{lstlisting}

Design notes worth being able to defend in review:
\begin{itemize}
    \item \texttt{S3KeySensor} with \texttt{poke\_interval=300} trades freshness against S3 LIST request cost; \texttt{timeout=3600} bounds the wait so a missing file fails visibly rather than hanging the run forever.
    \item The alert task uses \texttt{trigger\_rule="one\_failed"} so it runs precisely on the failure path---without it, the alert would be skipped when its upstream is skipped.
    \item The \texttt{\{\{ ds \}\}} template makes runs partition-aligned and rerunnable for any logical date---idempotency at the orchestration layer, matching Chapter~10's job-layer discipline.
\end{itemize}

\subsection*{Verification}
\begin{itemize}
    \item Uploading a matching file triggers a run that completes all four tasks.
    \item A run with a deliberately broken Glue job ends with the SNS alert published and the Airflow UI showing the failure path.
    \item A run with no file fails after one hour with a clear sensor timeout, not an indefinite hang.
    \item Rerunning a past logical date reprocesses exactly that date's partition (job-layer bookmarks honored).
\end{itemize}

\section{Friday Use Case: Self-Healing Banking Reconciliation Across Fifteen Systems}
\index{use case!banking reconciliation orchestration}
A bank reconciles positions daily across fifteen systems---core banking, card processor, SWIFT gateway, general ledger, and eleven more. Files arrive at different times with different formats; the close cannot complete until all fifteen reconcile; and yesterday's manual rerun procedure takes an engineer four hours. Design the orchestration.

\subsection{Architecture}
\begin{figure}[H]
    \centering
    \resizebox{\textwidth}{!}{%
    \begin{tikzpicture}[
        box/.style={rectangle, rounded corners, draw=primary, thick, fill=primary!6, align=center, minimum width=2.9cm, minimum height=1.0cm, font=\small\bfseries},
        guard/.style={rectangle, rounded corners, draw=red!60!black, thick, fill=red!5, align=center, minimum width=2.6cm, minimum height=0.9cm, font=\footnotesize\bfseries},
        arr/.style={-{Stealth[length=2.5mm]}, draw=primary, thick}
    ]
        \node[box] (eb) at (0,0) {EventBridge\\file arrivals};
        \node[box] (sf) at (3.4,0) {Step Functions\\per-system ingest};
        \node[box] (s3) at (6.8,0) {S3 staging\\per system/date};
        \node[box] (mwaa) at (10.2,0) {MWAA daily\\recon DAG};
        \node[box] (rs) at (13.2,0) {Redshift\\recon marts};
        \node[guard] (sns) at (10.2,2.0) {SNS alerts\\+ runbooks};
        \node[guard] (dq) at (6.8,-2.0) {Data-quality gate\\(DataBrew/DQ)};
        \draw[arr] (eb) -- (sf);
        \draw[arr] (sf) -- (s3);
        \draw[arr] (s3) -- (mwaa);
        \draw[arr] (mwaa) -- (rs);
        \draw[arr, dashed] (sns) -- (mwaa);
        \draw[arr, dashed] (dq) -- (s3);
    \end{tikzpicture}%
    }
    \caption{Two-tier orchestration: Step Functions handles per-arrival, per-system ingestion; MWAA owns the cross-system daily reconciliation with dependencies, retries, and alerting.}
    \label{fig:ch12-banking}
\end{figure}

\subsection{Design Decisions}
\textbf{Two tiers, each at its strength.} Step Functions runs fifteen per-system ingestion state machines, triggered by EventBridge on each file's arrival---event-driven, high-rate, AWS-native. MWAA owns the daily reconciliation DAG: one task per system reading that system's staged partition, a barrier task requiring all fifteen, the cross-system match, and the exception report. Step Functions is wrong for the cross-day dependency logic; Airflow is wrong for fifteen independent arrival streams. The boundary is the staged S3 partition.

\textbf{Self-healing by construction.} Every task is idempotent and partition-scoped (\texttt{system/date}), so any rerun is safe. Retries with backoff absorb transient failures. A sensor per system replaces human polling, with a timeout that escalates to the on-call with the system name and expected SLA. The ``manual rerun procedure'' becomes ``click \emph{Clear} on the failed task''---the four-hour incident becomes a five-minute one.

\textbf{Blast-radius controls.} Each system's task is isolated: a broken card-processor file fails that branch and alerts, while the other fourteen stage normally and the barrier reports exactly which system blocks the close. Partial close with flagged exceptions is a business decision; silent partial data is not an option.

\begin{pitfall}
\textbf{A single mega-DAG that recomputes everything on any failure.} If the failure of one of fifteen inputs forces reprocessing all fifteen, every incident is maximal. Partition-scoped, idempotent tasks with a barrier are what make reruns cheap enough to be the recovery strategy.
\end{pitfall}

\section{Interview Q\&A}
\subsection{Concepts and Trade-offs}
\begin{enumerate}
    \item \textbf{Q: What does \texttt{catchup=False} prevent when you deploy a DAG?}\\
    A: Backfill storms. With catchup enabled and a past \texttt{start\_date}, the scheduler immediately launches every missed historical run on deploy---often against production systems that did not expect it.
    \item \textbf{Q: When are Step Functions a better orchestrator than Airflow?}\\
    A: For event-driven, high-rate workflows composed of AWS service calls---per-file ingestion, per-order processing---where per-transition billing and native service integrations beat running an always-on scheduler.
    \item \textbf{Q: What do retries with backoff protect a pipeline against?}\\
    A: Transient failures---throttling, brief network or consistency hiccups---converting them into eventual success instead of pages. They do not fix deterministic failures; those need the failure path.
    \item \textbf{Q: Why tune an \texttt{S3KeySensor}'s \texttt{poke\_interval}?}\\
    A: Each poll is S3 LIST requests billed per thousand. A 300-second interval for a daily file balances freshness against request cost; a 5-second interval buys nothing and pays continuously.
    \item \textbf{Q: When should data move through S3 instead of XComs?}\\
    A: Always, for datasets. XComs live in the metadata database and suit small signals---keys, counts, statuses. Step Functions enforces the same rule with its 256~KB payload limit.
    \item \textbf{Q: Why does the alert task need \texttt{trigger\_rule="one\_failed"}?}\\
    A: Default trigger rules skip a task when its upstream is skipped. Without the explicit rule, the alert itself would be skipped on exactly the runs it exists for.
\end{enumerate}

\section{Further Reading}
The Airflow documentation is the reference for DAG semantics, scheduling, sensors, and trigger rules \cite{apacheAirflow}; the MWAA User Guide covers environment classes, worker scaling, and the S3 deployment model \cite{awsMWAADocs}; and the Step Functions Developer Guide covers ASL, service integrations, and Standard versus Express workflows \cite{awsStepFunctionsDocs}. The EventBridge guide covers the arrival-trigger half of the Friday architecture \cite{awsEventBridgeDocs}.

\section*{Key Takeaways}
\begin{itemize}
    \item Orchestration is dependency management, failure policy, and observability---computation belongs to the engines.
    \item Airflow's sharp edges are scheduling semantics: catchup, start dates, trigger rules, and retry policy.
    \item MWAA makes Airflow managed, not magic: execution-role least privilege and S3-based deployment are the operational core.
    \item Step Functions excels at event-driven AWS-native flows; its 256~KB payload limit enforces the S3-for-data rule.
    \item Sensors replace human polling; idempotent, partition-scoped tasks replace manual rerun procedures.
    \item Two-tier orchestration---native state machines for arrival, a DAG platform for cross-system logic---serves the 15-system reconciliation better than either tool alone.
\end{itemize}

\section{Exercises}
\begin{enumerate}
    \item \textbf{(Conceptual)} Explain the difference between a DAG, a DAG run, and a task instance, and which one a ``rerun for yesterday'' operates on.
    \item \textbf{(Conceptual)} Why is an always-on MWAA environment the wrong economic shape for a workload of pure per-arrival event reactions?
    \item \textbf{(Hands-on)} Deploy the Thursday DAG, then add a second sensor for a \texttt{control/} file and make the Glue task require both arrivals.
    \item \textbf{(Hands-on)} Force each failure mode---missing file, Glue failure, Athena failure---and record the alert behavior for each.
    \item \textbf{(Hands-on)} Reimplement the DAG's happy path as a Step Functions state machine and compare the operational evidence each tool gives you.
    \item \textbf{(Design)} Design the barrier task for the fifteen-system recon: how does it report exactly which systems are missing, and what happens at the close deadline?
    \item \textbf{(Design)} Write the MWAA execution-role policy for the Thursday DAG as least privilege; justify each statement.
    \item \textbf{(Architecture)} Extend the Friday architecture with a cross-region failover posture: what must exist in the second region, and what state is hard to replicate?
\end{enumerate}

\section{Milestone 3 Capstone: Automated Batch ELT Pipeline}\label{sec:ch12-capstone}
\index{capstone project}\index{Milestone 3 capstone}\index{Step Functions!capstone}

\begin{capstonebox}
\textbf{Mission.} Build an end-to-end automated batch ELT pipeline: an EventBridge rule fires on S3 upload, a Step Functions state machine runs a DataBrew data-quality check that gates an EMR Serverless Spark feature job, and the curated output loads into Redshift---with MWAA handling cross-day scheduling dependencies.

\textbf{Estimated time.} 6--8 hours in a sandbox AWS account.

\textbf{What you\textquotesingle{}ll practice.}
\begin{itemize}
    \item Event-driven triggering: S3 notifications to EventBridge to Step Functions.
    \item Data-quality gating inside a state machine with a Choice state.
    \item EMR Serverless job submission from orchestration.
    \item Failure routing to SNS with runbook-linked alerts.
\end{itemize}
\end{capstonebox}

\subsection*{Scenario}
The bank's analytics team receives daily feature extracts from an upstream vendor. Files arrive irregularly; quality varies; and the ML feature pipeline must never consume a failed-quality file. Automate arrival-to-Redshift with a hard quality gate and full failure visibility.

\subsection*{Requirements}
\textbf{Functional requirements.}
\begin{itemize}
    \item S3 event notifications to EventBridge; a rule matching the raw prefix starts the state machine with the object key in the payload.
    \item The state machine runs a DataBrew (or Glue Data Quality) check; a Choice state halts and alerts on failure.
    \item On pass, an EMR Serverless Spark job computes features; output is COPYed into Redshift.
    \item MWAA orchestrates the cross-day dependencies (yesterday's recon must finish before today's features publish).
    \item Failures at any state publish to SNS with the execution ARN.
\end{itemize}

\textbf{Non-functional requirements.}
\begin{itemize}
    \item End-to-end runs from a clean environment per the README: upload a file, pipeline completes.
    \item At least three automated data-quality checks with a failing-case demonstration that halts the load.
    \item CloudWatch alarms on \texttt{ExecutionsFailed} and on MWAA task failures.
    \item A monthly cost estimate (EMR Serverless vCPU-hours, Glue DPUs, Step Functions transitions) with two optimization decisions documented.
    \item Task-scoped IAM roles per pipeline stage; secrets in Secrets Manager; no hardcoded credentials anywhere.
\end{itemize}

\subsection*{Reference Architecture}
\begin{figure}[H]
    \centering
    \resizebox{\textwidth}{!}{%
    \begin{tikzpicture}[
        node distance=0.55cm and 1.0cm,
        box/.style={rectangle, rounded corners, draw=orange!75!black, thick, fill=orange!10, align=center, font=\small, minimum height=0.9cm, inner sep=4pt},
        side/.style={rectangle, rounded corners, draw=brown!70!black, thick, fill=brown!10, align=center, font=\footnotesize, inner sep=3pt},
        arr/.style={-{Stealth}, thick}
    ]
    \node[box] (s3) {S3 upload\\(raw file)};
    \node[box, right=of s3] (eb) {EventBridge\\rule};
    \node[box, right=of eb] (sf) {Step Functions\\state machine};
    \node[box, right=of sf, yshift=1.0cm] (dq) {Glue DataBrew\\DQ check};
    \node[box, right=of sf, yshift=-1.0cm] (emr) {EMR Serverless\\Spark features};
    \node[box, right=of emr, yshift=1.0cm] (rs) {Redshift\\(curated)};
    \node[side, below=0.6cm of sf] (mwaa) {MWAA (Airflow) cross-DAG deps};
    \draw[arr] (s3) -- (eb);
    \draw[arr] (eb) -- (sf);
    \draw[arr] (sf) -- (dq);
    \draw[arr] (sf) -- (emr);
    \draw[arr] (dq) -- (rs);
    \draw[arr] (emr) -- (rs);
    \draw[arr, dashed] (mwaa) -- (sf);
    \end{tikzpicture}%
    }
    \caption{Milestone 3 reference architecture: automated batch ELT pipeline with a data-quality gate.}
    \label{fig:ch12-capstone-arch}
\end{figure}

\subsection*{Implementation Steps}
\begin{enumerate}
    \item Enable S3 event notifications to EventBridge and create the prefix-matched rule targeting the state machine.
    \item Define the state machine: DataBrew job Task, the Choice gate, EMR Serverless Task, Redshift COPY Task, SNS failure catches.
    \item Create the EMR Serverless application and feature job; pass S3 keys---never data---between states (the 256~KB rule).
    \item Build the MWAA DAG that sequences the cross-day dependency.
    \item Wire the \texttt{ExecutionsFailed} alarm and MWAA failure alerts to SNS with runbook links.
    \item Write the cost model and the security review; tag resources for cleanup.
\end{enumerate}

The hardest state is the quality gate; the ASL Choice state is small but load-bearing:

\begin{lstlisting}[language=json,caption={The data-quality gate as a Choice state.},label={lst:ch12-choice}]
"CheckDataQuality": {
  "Type": "Choice",
  "Choices": [{
    "Variable": "$.dq_status",
    "StringEquals": "FAILED",
    "Next": "AlertAndHalt"
  }],
  "Default": "RunSparkFeatures"
}
\end{lstlisting}

\subsection*{Deliverables}
\begin{itemize}
    \item The state-machine definition (ASL), the DAG, and the EMR job code, all in version control.
    \item Evidence of the happy path (upload to Redshift rows) and of the DQ gate halting a bad file with an SNS alert.
    \item Alarm configurations and one triggered-alarm record.
    \item The cost estimate with two documented optimizations and the IAM/security review.
\end{itemize}

\subsection*{Acceptance Criteria}
\begin{itemize}
    \item A file upload triggers the pipeline with no manual step; the object key flows through as payload.
    \item A quality failure halts the load before any Spark or Redshift work and pages with context.
    \item No state passes bulk data through payloads; all data moves through S3 keys.
    \item Each pipeline stage runs under its own least-privilege role.
    \item The MWAA dependency demonstrably blocks today's publish until yesterday's recon succeeds.
\end{itemize}

\subsection*{Stretch Goals}
\begin{itemize}
    \item Add a Map state to fan out per-file processing when a manifest lists multiple files.
    \item Convert the DataBrew check to Glue Data Quality rules and compare rule expressiveness.
    \item Add Step Functions redrive policies and test recovery from a throttled EMR submission.
    \item Emit pipeline metrics (arrival-to-curated latency) to CloudWatch and dashboard the trend.
\end{itemize}
